% ============================================================
%  INFORME – !Blood: Aplicación Móvil para la Donación de Sangre
%  Formato: LaTeX  |  Compilar con pdflatex o Overleaf
% ============================================================
\documentclass[12pt, a4paper]{article}

% ── Paquetes ──────────────────────────────────────────────
\usepackage[utf8]{inputenc}
\usepackage[T1]{fontenc}
\usepackage[spanish]{babel}
\usepackage{geometry}
\geometry{left=2.5cm, right=2.5cm, top=2.5cm, bottom=2.5cm}
\usepackage{setspace}
\onehalfspacing
\usepackage{parskip}
\usepackage{titlesec}
\usepackage{enumitem}
\usepackage{graphicx}
\usepackage{xcolor}
\usepackage{hyperref}
\usepackage{fancyhdr}
\usepackage{booktabs}
\usepackage{array}
\usepackage{longtable}
\usepackage{caption}

% ── Colores institucionales ──────────────────────────────
\definecolor{bloodred}{HTML}{E53935}
\definecolor{darkgray}{HTML}{333333}
\definecolor{lightgray}{HTML}{F5F5F5}

% ── Estilos de hipervínculos ─────────────────────────────
\hypersetup{
    colorlinks=true,
    linkcolor=bloodred,
    urlcolor=bloodred,
    citecolor=bloodred
}

% ── Estilos de secciones ─────────────────────────────────
\titleformat{\section}
  {\Large\bfseries\color{bloodred}}{\thesection.}{0.5em}{}
\titleformat{\subsection}
  {\large\bfseries\color{darkgray}}{\thesubsection.}{0.5em}{}
\titleformat{\subsubsection}
  {\normalsize\bfseries\color{darkgray}}{\thesubsubsection.}{0.5em}{}

% ── Encabezado y pie ─────────────────────────────────────
\pagestyle{fancy}
\fancyhf{}
\fancyhead[L]{\small\textcolor{bloodred}{!Blood}}
\fancyhead[R]{\small Informe de Diseño}
\fancyfoot[C]{\thepage}
\renewcommand{\headrulewidth}{0.4pt}

% ============================================================
\begin{document}

% ── Portada ───────────────────────────────────────────────
\begin{titlepage}
\centering
\vspace*{3cm}

{\Huge\bfseries\textcolor{bloodred}{!Blood}\par}
\vspace{0.5cm}
{\Large Aplicación Móvil para la Donación de Sangre\par}
\vspace{1cm}
{\large\textit{Dona vida, recibe vida}\par}

\vspace{2cm}

{\Large\bfseries Informe de Diseño y Especificación Funcional\par}

\vspace{3cm}

{\large
Jesús Muñoz C.\\[0.3cm]
}

\vspace{2cm}

{\large Marzo 2026\par}

\vfill
\end{titlepage}

% ── Tabla de contenido ───────────────────────────────────
\newpage
\tableofcontents
\newpage

% ============================================================
% SECCIÓN 1: PLANTEAMIENTO DEL PROBLEMA
% ============================================================
\section{Planteamiento del Problema}

\subsection{Contexto}

La donación de sangre constituye un pilar fundamental de los sistemas de salud a escala mundial. De acuerdo con la Organización Mundial de la Salud (OMS), se requieren al menos entre 10 y 20 donaciones por cada 1\,000 habitantes para garantizar un suministro adecuado; sin embargo, numerosos países de América Latina, incluido Colombia, se encuentran significativamente por debajo de ese umbral. En el departamento del Cauca, y particularmente en Popayán, la situación se agrava por la dispersión geográfica y por la ausencia de mecanismos ágiles que conecten a potenciales donantes con quienes requieren transfusiones de forma urgente.

\subsection{Problemática Identificada}

A partir del análisis del entorno se identificaron los siguientes problemas centrales:

\begin{enumerate}[label=\textbf{\arabic*.}, leftmargin=1.5cm]
    \item \textbf{Información dispersa y desactualizada.} Las convocatorias de donación circulan por redes sociales, volantes o comunicados institucionales sin un canal unificado. El ciudadano desconoce cuál fuente es oficial, cuándo caduca la necesidad o cuáles centros de salud están habilitados.

    \item \textbf{Desconexión entre la oferta y la demanda.} Cuando un paciente necesita una transfusión, su familia recurre a cadenas de mensajes informales (WhatsApp, Facebook) sin posibilidad de filtrar por compatibilidad de grupo sanguíneo ni por cercanía geográfica, lo que dilata los tiempos de respuesta.

    \item \textbf{Desinformación y mitos.} Creencias populares ---como que donar sangre debilita al organismo, afecta el peso corporal o compromete la vida sexual--- disuaden a potenciales donantes que, en condiciones de información correcta, estarían dispuestos a participar.

    \item \textbf{Falta de seguimiento y motivación.} El donante eventual no recibe retroalimentación sobre el impacto de su donación, carece de un historial accesible y no dispone de recordatorios sobre cuándo puede volver a donar, lo que reduce la tasa de donación recurrente.

    \item \textbf{Ausencia de herramientas para profesionales de salud.} Los centros de salud y hospitales no cuentan con un medio digital eficiente para publicar campañas de donación, gestionar cupos ni comunicar requisitos de manera estandarizada a la comunidad.

    \item \textbf{Comunicación insegura y no moderada.} Los canales informales de comunicación entre solicitantes y donantes no garantizan la privacidad de los datos ni limitan las interacciones a fines estrictamente relacionados con la donación.
\end{enumerate}

\subsection{Justificación}

Frente a la problemática descrita, surge la necesidad de una herramienta tecnológica que centralice el proceso de donación de sangre, conecte de manera inteligente a donantes y solicitantes con base en compatibilidad sanguínea y proximidad geográfica, y proporcione un entorno seguro y moderado de comunicación. Una solución móvil ---dado que el teléfono inteligente es el dispositivo de uso predominante en la población objetivo--- puede reducir las barreras de acceso, incrementar la tasa de donación voluntaria y, como consecuencia directa, contribuir a salvar vidas.

% ============================================================
% SECCIÓN 2: DESCRIPCIÓN DE LA HERRAMIENTA DE SOFTWARE
% ============================================================
\section{Descripción de la Herramienta de Software}

\subsection{Definición General}

!Blood es una aplicación móvil orientada a la donación de sangre que permite a cualquier persona integrarse a una comunidad de apoyo, con la posibilidad de donar o solicitar sangre según la necesidad. La herramienta centraliza solicitudes, campañas y comunicación entre usuarios, priorizando la compatibilidad por tipo de sangre y la cercanía geográfica, para facilitar respuestas oportunas.

La aplicación contempla dos perfiles de usuario diferenciados ---ciudadano y profesional de salud--- y estructura su funcionamiento en torno a módulos claramente definidos que cubren desde el registro y verificación de identidad hasta la gestión de campañas institucionales.

\subsection{Arquitectura Funcional}

A continuación se describen en detalle los módulos funcionales que componen la aplicación, organizados según el flujo natural de interacción del usuario.

% ── 2.2.1 Pantalla de Bienvenida ─────────────────────────
\subsubsection{Pantalla de Bienvenida (Splash)}

La pantalla de bienvenida constituye el primer punto de contacto del usuario con la aplicación. Presenta la identidad visual de !Blood ---logotipo con la gota de sangre y el eslogan ``Dona vida, recibe vida''--- acompañada de un mensaje de comunidad: \textit{``Juntos creamos una comunidad que salva vidas''}. Desde esta pantalla, el usuario accede a dos acciones principales:

\begin{itemize}
    \item \textbf{Iniciar Sesión:} redirige a la pantalla de autenticación para usuarios ya registrados.
    \item \textbf{Registrarme:} inicia el flujo de creación de cuenta en tres pasos.
\end{itemize}

% ── 2.2.2 Inicio de Sesión ───────────────────────────────
\subsubsection{Inicio de Sesión (Login)}

El módulo de inicio de sesión permite al usuario autenticarse mediante su correo electrónico o nombre de usuario y su contraseña. La interfaz incluye:

\begin{itemize}
    \item Campo de usuario o correo electrónico con icono contextual.
    \item Campo de contraseña con opción de alternar la visibilidad del texto ingresado (botón de ojo).
    \item Enlace de recuperación de contraseña (\textit{``¿Olvidaste tu contraseña?''}).
    \item Enlace de acceso rápido al registro para usuarios nuevos.
\end{itemize}

El diseño prioriza la simplicidad: formulario limpio, botones claros y retroalimentación visual inmediata.

% ── 2.2.3 Registro de Usuario ────────────────────────────
\subsubsection{Registro de Usuario}

El proceso de registro se estructura en tres pasos secuenciales, con un indicador de progreso visible que informa al usuario en qué etapa se encuentra.

\paragraph{Paso 1: Datos Personales.}
En esta primera etapa, el usuario proporciona:

\begin{itemize}
    \item \textbf{Tipo de usuario:} selección entre \textit{Ciudadano} o \textit{Profesional de salud}. Esta distinción habilita funcionalidades diferenciadas (por ejemplo, la creación de campañas está reservada al perfil profesional).
    \item \textbf{Nombre completo.}
    \item \textbf{Fecha de nacimiento.}
    \item \textbf{Tipo de sangre:} selección entre los ocho grupos sanguíneos (A+, A-, B+, B-, AB+, AB-, O+, O-).
    \item \textbf{Cédula de ciudadanía.}
    \item \textbf{Teléfono de contacto.}
    \item \textbf{Correo electrónico.}
\end{itemize}

\paragraph{Paso 2: Verificación.}
La segunda etapa tiene como objetivo validar la identidad del usuario y registrar su ubicación:

\begin{itemize}
    \item \textbf{Fotografía de la cédula de ciudadanía:} el usuario carga una imagen legible de su documento de identidad.
    \item \textbf{Foto de perfil:} se solicita una fotografía con fondo blanco y rostro visible, para facilitar la identificación entre usuarios.
    \item \textbf{Ciudad de residencia.}
    \item \textbf{Dirección:} campo de texto complementado con un botón de geolocalización automática que permite capturar la ubicación actual del dispositivo.
\end{itemize}

\paragraph{Paso 3: Configuración de Cuenta.}
En la última etapa del registro se configuran las credenciales de acceso:

\begin{itemize}
    \item \textbf{Nombre de usuario:} identificador único dentro de la plataforma.
    \item \textbf{Contraseña:} con indicador visual de fortaleza (baja, media, alta) que orienta al usuario a crear una contraseña segura.
    \item \textbf{Confirmación de contraseña.}
    \item \textbf{Aceptación de términos y condiciones} y de la política de privacidad.
    \item \textbf{Compromiso de donación solidaria:} el usuario acepta explícitamente que toda donación canalizada a través de la plataforma es voluntaria y sin fines de lucro.
\end{itemize}

% ── 2.2.4 Pantalla Principal ─────────────────────────────
\subsubsection{Pantalla Principal (Home)}

La pantalla principal es el centro de operaciones del usuario tras iniciar sesión. Está compuesta por los siguientes elementos:

\paragraph{Encabezado personalizado.}
Muestra el avatar, el saludo personalizado con el nombre del usuario y un botón de acceso directo a las notificaciones, con indicador numérico de notificaciones pendientes.

\paragraph{Tarjeta de sangre del usuario.}
Panel destacado que presenta el tipo de sangre del usuario junto con sus estadísticas de participación: número de donaciones realizadas y vidas ayudadas. Este elemento refuerza la motivación y el sentido de impacto.

\paragraph{Acciones rápidas.}
Cuadrícula de cuatro accesos directos a las funcionalidades principales:

\begin{itemize}
    \item \textbf{Solicitar sangre:} abre el formulario de creación de solicitud.
    \item \textbf{Ver solicitudes:} accede al listado de solicitudes activas en la zona.
    \item \textbf{Campañas:} muestra las campañas de donación programadas.
    \item \textbf{Mensajes:} dirige al módulo de chat.
\end{itemize}

\paragraph{Solicitudes cercanas.}
Listado resumido de las solicitudes de sangre activas más próximas al usuario, mostrando para cada una: nombre del solicitante, tipo de sangre requerido, distancia en kilómetros y tiempo transcurrido desde la publicación. Cada elemento es interactivo y conduce al detalle de la solicitud.

\paragraph{Campañas próximas.}
Vista previa de las campañas de donación programadas en la zona, con información de la institución organizadora, fecha, hora y distancia. Este componente permite al usuario conocer las jornadas de donación sin necesidad de navegar a una sección separada.

\paragraph{Estadísticas de comunidad.}
Indicadores globales que muestran el total de donantes activos y vidas salvadas en la plataforma, reforzando el sentido de pertenencia a una comunidad con impacto real.

\paragraph{Barra de navegación inferior.}
Barra fija con cinco elementos: Inicio, Explorar, botón central prominente de creación rápida (+), Chat y Perfil. El botón central se destaca visualmente como un botón de acción flotante (FAB) para facilitar la creación de solicitudes.

% ── 2.2.5 Creación de Solicitud ──────────────────────────
\subsubsection{Creación de Solicitud de Sangre}

Este módulo permite a cualquier usuario publicar una solicitud de sangre. El formulario recoge la siguiente información:

\begin{itemize}
    \item \textbf{Tipo de sangre necesario:} selector visual con los ocho grupos sanguíneos dispuestos en una cuadrícula de fácil selección.
    \item \textbf{Unidades necesarias:} control numérico con botones de incremento y decremento (rango de 1 a 10 unidades).
    \item \textbf{Nivel de urgencia:} tres opciones claramente diferenciadas por color e iconografía:
    \begin{itemize}
        \item \textit{Bajo:} puede esperar días.
        \item \textit{Medio:} se necesita en las próximas horas.
        \item \textit{Urgente:} se necesita lo antes posible.
    \end{itemize}
    \item \textbf{Centro de salud:} selección del centro médico donde se realizará la donación. Un aviso informativo aclara que la donación se efectúa exclusivamente en dicho centro.
    \item \textbf{Mensaje para donantes:} campo de texto opcional (máximo 200 caracteres) donde el solicitante puede describir brevemente su situación y apelar a la solidaridad de la comunidad.
\end{itemize}

Un banner informativo en la parte superior advierte: \textit{``Tu solicitud será visible para donantes compatibles en tu zona''}, garantizando transparencia sobre el alcance de la publicación.

\paragraph{Confirmación de publicación.}
Tras enviar la solicitud, se muestra una pantalla de éxito que informa al usuario:
\begin{itemize}
    \item El número de donantes compatibles cercanos que fueron notificados.
    \item Que recibirá notificaciones cuando algún donante acepte.
    \item Accesos directos para volver al inicio o consultar sus solicitudes activas.
\end{itemize}

% ── 2.2.6 Listado de Solicitudes ─────────────────────────
\subsubsection{Exploración de Solicitudes}

La pantalla de solicitudes presenta un listado completo de las solicitudes de sangre activas en la zona. Incluye las siguientes herramientas de filtrado y búsqueda:

\begin{itemize}
    \item \textbf{Barra de búsqueda:} permite buscar por nombre del solicitante o tipo de sangre.
    \item \textbf{Filtros rápidos (chips):} botones de filtrado por categorías: Todas, Urgentes, Compatible con el tipo de sangre del usuario, y Cercanas.
\end{itemize}

Cada tarjeta de solicitud muestra:
\begin{itemize}
    \item Indicador visual de nivel de urgencia (color y etiqueta).
    \item Tiempo transcurrido desde la publicación.
    \item Fotografía y nombre del solicitante.
    \item Tipo de sangre requerido y cantidad de unidades.
    \item Centro de salud asociado.
    \item Distancia desde la ubicación del usuario.
    \item Mensaje breve del solicitante.
\end{itemize}

% ── 2.2.7 Detalle de Solicitud ───────────────────────────
\subsubsection{Detalle de Solicitud}

Al seleccionar una solicitud, el usuario accede a una vista detallada que comprende:

\begin{itemize}
    \item \textbf{Encabezado visual:} indicador de urgencia, tipo de sangre en formato grande y cantidad de unidades necesarias.
    \item \textbf{Información del solicitante:} fotografía, nombre, antigüedad como miembro de la comunidad y sello de verificación de identidad.
    \item \textbf{Mensaje completo:} descripción detallada de la situación del solicitante.
    \item \textbf{Centro de salud:} nombre, dirección, distancia y botón para visualizar en mapa.
    \item \textbf{Información adicional:} tiempo transcurrido y estado de donantes aceptados respecto al total requerido (por ejemplo, ``1 de 3'').
    \item \textbf{Aviso de compatibilidad:} indicador automático que informa al usuario si su tipo de sangre es compatible con el requerido.
    \item \textbf{Botón de acción:} \textit{``Quiero donar''}, que conduce al flujo de confirmación.
    \item \textbf{Nota informativa:} aclara que al aceptar se habilita un chat interno para coordinar y que la donación se realiza únicamente en el centro de salud indicado.
\end{itemize}

% ── 2.2.8 Confirmación de Donación ───────────────────────
\subsubsection{Confirmación de Donación}

Antes de confirmar su intención de donar, el usuario accede a una pantalla que consolida:

\begin{itemize}
    \item \textbf{Resumen de la solicitud:} solicitante, tipo de sangre y centro de salud.
    \item \textbf{Compromiso del donante:} lista de compromisos que el usuario acepta:
    \begin{itemize}
        \item Acudir al centro de salud indicado.
        \item La donación es voluntaria y sin fines de lucro.
        \item La comunicación se limitará a la coordinación de la donación.
    \end{itemize}
    \item \textbf{Confirmación de salud:} lista de verificación (checklist) donde el usuario declara:
    \begin{itemize}
        \item No haber donado sangre en los últimos 3 meses.
        \item No tener ninguna enfermedad infectocontagiosa.
        \item Tener entre 18 y 65 años.
        \item Pesar más de 50\,kg.
    \end{itemize}
\end{itemize}

\paragraph{Pantalla de agradecimiento.}
Al confirmar, se muestra una pantalla de reconocimiento con el mensaje \textit{``¡Gracias por donar!''}, indicando que el donante y el solicitante ya pueden coordinarse vía chat y recordando el centro de salud al que debe acudir.

% ── 2.2.9 Sistema de Mensajería ──────────────────────────
\subsubsection{Sistema de Mensajería (Chat)}

El módulo de mensajería permite la comunicación directa entre donante y solicitante, exclusivamente en el contexto de una donación aceptada.

\paragraph{Lista de conversaciones.}
La pantalla de mensajes presenta todas las conversaciones activas del usuario, mostrando para cada una:
\begin{itemize}
    \item Fotografía del interlocutor con indicador de estado en línea.
    \item Nombre y última hora de actividad.
    \item Vista previa del último mensaje.
    \item Indicador de mensajes no leídos.
    \item Contexto de la conversación: tipo de sangre de la solicitud asociada y estado (activa o completada).
\end{itemize}

\paragraph{Detalle de conversación.}
La vista de chat individual incluye:
\begin{itemize}
    \item Encabezado con datos del interlocutor, solicitud de sangre asociada y centro de salud.
    \item Área de mensajes con burbujas diferenciadas para mensajes enviados y recibidos, con marcas de tiempo e indicadores de lectura (doble check).
    \item Separadores de fecha para organizar cronológicamente la conversación.
    \item \textbf{Aviso de moderación por IA:} banner visible que informa: \textit{``Chat protegido por IA -- Solo mensajes relacionados con la donación''}. Este mecanismo asegura que la comunicación se mantenga dentro del propósito de la plataforma, filtrando contenido inadecuado o no relacionado.
    \item Barra de entrada de mensaje con campo de texto y botón de envío.
\end{itemize}

% ── 2.2.10 Perfil de Usuario ─────────────────────────────
\subsubsection{Perfil de Usuario}

La pantalla de perfil concentra la información personal y las opciones de configuración del usuario:

\begin{itemize}
    \item \textbf{Cabecera visual:} imagen de portada, foto de perfil con opción de edición (botón de cámara), nombre completo y nombre de usuario.
    \item \textbf{Insignias:} sellos de \textit{Verificado} y \textit{Donante activo} que reconocen el estatus del usuario.
    \item \textbf{Tarjeta de estadísticas:} tipo de sangre del usuario junto con tres métricas: donaciones realizadas, vidas ayudadas y solicitudes creadas.
    \item \textbf{Información personal:} correo electrónico, teléfono, ubicación y antigüedad como miembro.
    \item \textbf{Opciones de configuración:}
    \begin{itemize}
        \item Editar perfil.
        \item Configuración de notificaciones.
        \item Configuración de privacidad.
        \item Centro de ayuda.
        \item Cerrar sesión.
    \end{itemize}
\end{itemize}

% ── 2.2.11 Notificaciones ────────────────────────────────
\subsubsection{Sistema de Notificaciones}

El módulo de notificaciones mantiene al usuario informado en tiempo real sobre las novedades relevantes. Cada notificación incluye un icono contextual, título, descripción y marca temporal. Los tipos de notificación identificados son:

\begin{itemize}
    \item \textbf{Solicitudes urgentes:} alerta cuando se publica una solicitud de sangre compatible y cercana.
    \item \textbf{Mensajes nuevos:} aviso de que un usuario ha enviado un mensaje en el chat.
    \item \textbf{Donaciones confirmadas:} confirmación de que una donación fue registrada exitosamente.
    \item \textbf{Reconocimientos:} mensajes motivacionales que celebran los logros del usuario (\textit{``Ya llevas 2 donaciones. Eres un héroe de la comunidad''}).
    \item \textbf{Estadísticas de comunidad:} actualizaciones sobre el crecimiento de la base de usuarios.
\end{itemize}

Las notificaciones se distinguen visualmente entre leídas y no leídas, y se ofrece un botón global para marcar todas como leídas.

% ── 2.2.12 Mis Solicitudes ───────────────────────────────
\subsubsection{Gestión de Solicitudes Propias}

El usuario puede consultar y administrar sus propias solicitudes de sangre mediante una interfaz organizada en pestañas: \textit{Activas}, \textit{Completadas} y \textit{Canceladas}.

Para cada solicitud activa se presenta:
\begin{itemize}
    \item Tipo de sangre y estado de la solicitud.
    \item Centro de salud y fecha de creación.
    \item Barra de progreso que indica la cantidad de donantes confirmados respecto al total requerido.
    \item Avatares de los donantes que han aceptado.
    \item Botones de acción: \textit{Ver detalles} y \textit{Cancelar}.
\end{itemize}

% ── 2.2.13 Campañas de Donación ──────────────────────────
\subsubsection{Campañas de Donación}

Las campañas de donación representan una funcionalidad exclusiva para usuarios con perfil de \textbf{profesional de salud}. Permiten a instituciones de salud organizar y promocionar jornadas de donación masivas.

\paragraph{Creación de campaña.}
El formulario de creación recoge:

\begin{itemize}
    \item \textbf{Nombre de la campaña.}
    \item \textbf{Centro de salud o institución organizadora.}
    \item \textbf{Dirección del evento:} con soporte de geolocalización automática.
    \item \textbf{Fecha, hora de inicio y hora de finalización.}
    \item \textbf{Cupos disponibles:} número máximo de donantes que se pueden recibir.
    \item \textbf{Tipos de sangre necesarios:} selección múltiple que permite especificar qué grupos sanguíneos se requieren, o seleccionar ``Todos''.
    \item \textbf{Descripción de la campaña:} campo de texto amplio (máximo 500 caracteres) para detallar los objetivos, beneficios para donantes y cualquier información relevante.
    \item \textbf{Requisitos para donar:} lista estandarizada de requisitos que incluye edad (18--65 años), peso mínimo (50\,kg), documento de identidad y tiempo mínimo desde la última donación (3 meses).
    \item \textbf{Teléfono de contacto:} para consultas adicionales.
\end{itemize}

Un banner informativo comunica: \textit{``Las campañas se mostrarán a usuarios cercanos al lugar del evento''}, asegurando que el organizador comprenda el alcance de la publicación.

\paragraph{Confirmación de publicación.}
Al publicar, se muestra una pantalla de éxito con:
\begin{itemize}
    \item Número de usuarios cercanos que serán notificados.
    \item Fecha y ubicación de la campaña como recordatorio.
    \item Acceso a la gestión de campañas del usuario.
\end{itemize}

\paragraph{Exploración de campañas.}
Todos los usuarios de la plataforma pueden explorar las campañas disponibles a través de un listado que incluye:
\begin{itemize}
    \item Barra de búsqueda.
    \item Filtros rápidos: Cercanas, Esta semana, Este mes.
    \item Tarjetas de campaña con: distintivo visual con fecha, nombre, institución organizadora, dirección, horario, distancia y cupos disponibles respecto al total.
\end{itemize}

\paragraph{Detalle de campaña.}
La vista detallada presenta de forma exhaustiva:
\begin{itemize}
    \item Encabezado visual con fecha destacada.
    \item Información del organizador con sello de verificación institucional.
    \item Descripción completa de la campaña.
    \item Mapa de ubicación con distancia y botón de indicaciones.
    \item Horario detallado.
    \item Tipos de sangre necesarios con indicación de los más urgentes.
    \item Requisitos para donar.
    \item Indicador de disponibilidad: cupos reservados vs.\ disponibles, con barra de progreso visual.
    \item Información de contacto.
    \item Botón de acción \textit{``Reservar mi cupo''} con nota informativa que aclara que la reserva no es obligatoria, pero ayuda a la organización logística.
\end{itemize}

% ── 2.2.14 Resumen de características ────────────────────
\subsection{Características Transversales}

\begin{itemize}
    \item \textbf{Compatibilidad inteligente:} el sistema evalúa automáticamente la compatibilidad entre el tipo de sangre del donante y el tipo requerido, notificando al usuario cuando puede ayudar.
    \item \textbf{Geolocalización:} todas las solicitudes y campañas muestran la distancia respecto a la ubicación del usuario, priorizando las más cercanas.
    \item \textbf{Verificación de identidad:} el registro requiere cédula de ciudadanía y fotografía, generando un sello de ``Verificado'' que otorga confianza dentro de la comunidad.
    \item \textbf{Moderación por inteligencia artificial:} el chat incorpora filtrado automático para garantizar que la comunicación se limite al contexto de la donación, protegiendo la privacidad y seguridad de los usuarios.
    \item \textbf{Modelo solidario:} la plataforma refuerza de forma explícita, tanto en el registro como en el flujo de donación, que las donaciones son voluntarias y sin fines de lucro.
    \item \textbf{Gamificación y reconocimiento:} estadísticas de donaciones, insignias y mensajes motivacionales fomentan la participación recurrente.
    \item \textbf{Interfaz inclusiva:} diseño responsive con tipografía clara (Inter), iconografía intuitiva (FontAwesome), campos de formulario con pistas contextuales y flujos cortos que minimizan la carga cognitiva.
\end{itemize}

% ============================================================
% SECCIÓN 3: PERFILES DE USUARIO
% ============================================================
\section{Perfiles de Usuario}

% ── 3.1 Donante Potencial ────────────────────────────────
\subsection{Perfil 1: Donante Potencial}

\subsubsection{Datos Demográficos}
Es una persona entre 18 y 65 años que vive en Popayán o sus alrededores. Cuenta con documento de identidad y suele movilizarse por estudio, trabajo o diligencias dentro de la ciudad. Su peso es mayor a 50\,kg.

\subsubsection{Objetivos y Necesidades}
Quiere ayudar donando sangre, pero necesita que le expliquen de forma sencilla si cumple los requisitos y qué debe hacer antes de ir. También necesita saber con claridad dónde puede donar, si hay una jornada cerca, y poder agendar una hora para no perder tiempo.

\subsubsection{Motivaciones}
Le motiva salvar vidas, apoyar a su comunidad y sentir que su acción tiene un impacto real. Responde mejor cuando entiende que lo necesitan y que el proceso es rápido y seguro.

\subsubsection{Frustraciones o Problemas}
Se frustra cuando la información está dispersa en páginas web o redes sociales y no sabe cuál es la fuente oficial o actualizada. Se desanima si llega a donar y le informan en el punto que no puede hacerlo por algo que nadie le explicó antes. También le generan dudas los mitos, por ejemplo creer que donar lo debilita, afecta su peso o su vida sexual, o que la sangre se vende.

\subsubsection{Habilidades Tecnológicas}
Usa el celular a diario para redes sociales, mensajería y trámites básicos. Puede registrarse en una aplicación si el proceso es corto y el lenguaje es claro.

\subsubsection{Contexto de Uso}
Normalmente revisa el teléfono en ratos libres. Si recibe una alerta de necesidad o una jornada de donación, decide rápido si puede ir, pero necesita saber el tiempo estimado y tener confirmación.

\subsubsection{Comportamientos y Hábitos Relevantes}
Si es hombre puede donar cada tres meses y si es mujer cada seis meses, por lo que necesita que la aplicación le recuerde cuándo vuelve a ser apto. Puede tener dudas sobre temas sensibles del requisito (por ejemplo, cambios de pareja sexual o antecedentes de hepatitis) y por eso requiere un cuestionario privado, respetuoso y fácil de entender.

% ── 3.2 Solicitante de Sangre ────────────────────────────
\subsection{Perfil 2: Solicitante de Sangre}

\subsubsection{Datos Demográficos}
Persona adulta (familiar, amigo o cuidador de un paciente) que reside en Popayán o se encuentra temporalmente en la ciudad por motivos de atención médica. Puede pertenecer a cualquier estrato socioeconómico y grado de formación académica.

\subsubsection{Objetivos y Necesidades}
Necesita conseguir donantes de un tipo de sangre específico en el menor tiempo posible para un paciente que se encuentra en un centro de salud. Requiere una herramienta que le permita publicar la necesidad, llegar a personas compatibles cercanas y coordinar la logística de la donación de manera clara y segura.

\subsubsection{Motivaciones}
La urgencia de la situación médica de un ser querido lo impulsa a actuar con rapidez. Valora cualquier medio que le permita ampliar su alcance más allá de su círculo social inmediato.

\subsubsection{Frustraciones o Problemas}
Las cadenas de mensajes en redes sociales o aplicaciones de mensajería no le permiten filtrar por compatibilidad ni por cercanía. No tiene forma de saber quién realmente puede o quiere donar. La incertidumbre y la espera generan angustia adicional en un contexto ya estresante.

\subsubsection{Habilidades Tecnológicas}
Similares al donante potencial; usa el celular para comunicaciones básicas. En una situación de urgencia, necesita que el proceso de publicación sea extremadamente rápido e intuitivo.

\subsubsection{Contexto de Uso}
Se encuentra en un centro de salud o cerca de él, con acceso a su teléfono móvil. Necesita publicar la solicitud, recibir respuestas y coordinar con donantes, todo desde la misma herramienta.

\subsubsection{Comportamientos y Hábitos Relevantes}
En situación de emergencia, es probable que su capacidad de atención y paciencia se vean reducidas. Por ello, la aplicación debe minimizar los pasos requeridos y ofrecer retroalimentación constante sobre el estado de su solicitud (donantes notificados, donantes aceptados, etc.).

% ── 3.3 Profesional de Salud ─────────────────────────────
\subsection{Perfil 3: Profesional de Salud}

\subsubsection{Datos Demográficos}
Profesional vinculado a una institución de salud (hospital, clínica, banco de sangre, universidad con facultad de medicina) en Popayán o sus alrededores. Puede ser médico, enfermero, coordinador de banco de sangre o personal administrativo autorizado.

\subsubsection{Objetivos y Necesidades}
Busca una herramienta que le permita crear y publicar campañas de donación para su institución, alcanzando a la mayor cantidad posible de potenciales donantes en la zona. Necesita gestionar cupos, comunicar requisitos y fechas, y hacer seguimiento de las reservas realizadas.

\subsubsection{Motivaciones}
Contribuir al abastecimiento de sangre en su institución y en la comunidad. Le motiva contar con un canal digital directo que complemente los métodos tradicionales de convocatoria.

\subsubsection{Frustraciones o Problemas}
Los medios tradicionales de convocatoria (volantes, radio, redes sociales institucionales) tienen alcance limitado y no permiten medir respuestas ni organizar la logística de donantes. Las campañas exitosas dependen de factores no controlables como la viralización espontánea de la información.

\subsubsection{Habilidades Tecnológicas}
Posee habilidades digitales medias a altas. Está acostumbrado a utilizar sistemas de información hospitalaria y aplicaciones profesionales.

\subsubsection{Contexto de Uso}
Planifica las campañas desde su lugar de trabajo, con tiempo suficiente para completar formularios detallados. Consulta el estado de las reservas y las estadísticas de la campaña de forma periódica.

\subsubsection{Comportamientos y Hábitos Relevantes}
Valora la formalidad y el respaldo institucional. Necesita que la plataforma le permita publicar bajo el nombre de su institución y que la campaña se muestre como un contenido verificado y confiable.

% ============================================================
% Cierre
% ============================================================
\newpage
\section*{Conclusión}
\addcontentsline{toc}{section}{Conclusión}

!Blood se concibe como una solución integral al problema de la desconexión entre la necesidad de sangre y la disposición a donar. Mediante un diseño centrado en el usuario, la aplicación reduce las barreras de información, simplifica la coordinación logística, protege la seguridad de los datos y refuerza el modelo de donación voluntaria y solidaria. Los tres perfiles de usuario identificados ---donante potencial, solicitante de sangre y profesional de salud--- representan los actores fundamentales del ecosistema de donación, y la herramienta está diseñada para satisfacer las necesidades específicas de cada uno a través de funcionalidades diferenciadas, accesibles e intuitivas.

\end{document}
